% 原文前言：印刷vii–xiii／PDF物理8–14.
\section*{前言}\phantomsection\label{sec:preface}\addcontentsline{toc}{section}{前言}
\subsection*{为什么写这本书？}
为什么还要写一本最优运输（optimal transport）的书？菲尔兹奖得主Cédric Villani的两本书还不够吗？还有Luigi Ambrosio、Nicola Gigli和Giuseppe Savaré撰写的梯度流“圣经”，其中同样包含许多关于最优运输的高级而一般的细节.

本书很松散地受到我2011、2012年在奥赛开设的课程，以及更早在格勒诺布尔暑期学校讲授的两门短期入门课程\cite{b274,b275}的启发，希望提出不同的视角，也部分面向不同的读者群体. 目前，将最优运输作为工具从事应用数学研究的群体正在不断扩大. 例如图像处理、经济学和演化偏微分方程，后者尤其涉及生物学或社会科学中的群体动力学建模，以及流体力学. 更一般地，在应用数学中，最优运输既是进行证明、估计和设计数值方法的技术工具，也是描述涉及距离、路径和成本的现象的建模工具.

从这一侧进入最优运输的读者最关心的问题包括：怎样数值计算最优映射与成本；问题有哪些不同表述，以便将其用于新的模型和目的；怎样以尽可能具体的方式得到连接运输与偏微分方程的主要结果；以及怎样运用该理论研究已有模型. 我们不能说已有著作没有回答这些问题，但很可能没有一本是以此为写作目的的.

C. Villani的第一本书\cite{b292}最接近我们的意图. 它广泛介绍这一主题及其应用，适合各种读者. 但我决定在这里讨论的一些主题并未在\cite{b292}中描述，例如本书\hyperref[sec:4]{第4章}的最小流问题. 而且，自2003年以来，这一理论已经取得巨大进展. ST Rachev和L Ruschednorf的两卷著作\cite{b257,b258}也不应被忘记，其中涉及许多概率与统计应用. 不过，它们的研究范围很快便与本书分开，因此我们不会展开\cite{b257}中的大部分应用和语言. 我们主要坚持确定性的框架与变分法的风格.

如果考察\cite{b292}之后的发展，至少应提到此后出版的两部优美著作. C. Villani的新书\cite{b293}将原先的介绍扩充为近千页的篇幅，增加的内容大多涉及几何问题，尤其是曲率概念. 流形上二次成本的最优运输成为核心工具，也是研究度量测度空间的出发点. 另一部参考书是L. Ambrosio、N. Gigli和G. Savaré的著作\cite{b15}，研究度量空间，尤其是概率测度的Wasserstein空间中的梯度流演化. 这一主题有许多有趣应用，例如热方程、Fokker--Planck方程与多孔介质方程，所发展的工具也非常有用. 但作者的主要视角是研究这些抽象空间中隐藏的微分结构；建模与应用数学大概并非他们首先关注的事情.

上述部分文献技术性很强，发展了能够适用于非常一般情形的强大抽象工具，但许多读者可能难以使用. 因而后来出现了一些较短的综述和更容易处理的讲义，例如\cite{b11}对\cite{b15}部分内容的简化，以及L. Ambrosio和N. Gigli撰写的简短入门\cite{b9}. 然而，仔细分析\cite{b9}便会发现，前半部分介绍带有某些变化的最优运输一般理论，余下部分讨论一般度量空间的梯度流，以及具有曲率下界的度量测度空间. 还应提到两篇很新的综述\cite{b63,b228}：前者在L. V. Kantorovich诞辰百年之际回顾该理论的主要成果，后者则对这一主题及其与几何和经济学的联系作总体介绍.

与此同时，许多活动凸显了人们对最优运输应用侧面日益增长的兴趣：出版物\footnote{ESAIM M2AN的“Optimal Transport in Applied Mathematics”专刊正在筹备；本书部分参考文献取自该专刊.}、学校、研讨会、研究项目等. 自2010年代以来，这个群体十分活跃，尤其在法国，也包括加拿大、意大利、奥地利、英国等地. 这些考虑使我觉得，撰写一本专门的书可能是个好主意；你现在手中的这本书，就是过去三年工作的成果.

\subsection*{这本书有什么特点？}
本书严格介绍最优运输理论及一些受到忽视的变体，并解释它与演化偏微分方程、图像处理和经济学中许多主题之间最重要的联系.

我尽可能避开最一般的框架，将注意力集中于欧氏情形；但若在Polish空间中陈述结果并不增加困难，反而有助于澄清图景，也会采用这一表述. 只要加上紧性假设能够简化叙述，又不至于过多降低结论的意义——当然是依照我的个人判断——我便借此略过一些困难. 许多时候，我也先从最简单的情形，例如具有紧性或连续性的情形讲起，然后再推广.

在可以选择时，我尽量采用更“具体”的证明，因为我认为读者更容易理解. 例如，Wasserstein空间中Lipschitz曲线的速度场存在性，在这里通过逼近证明；而\cite{b15}采用抽象泛函分析工具，关键是巧妙运用Hahn--Banach定理.

我没有试图穷尽所有可能的主题，而是把全书分为八章，大致各对应一讲较长的课程. 显然，我增加了许多通常一讲之内无法涉及的材料，但主题的选择与顺序确实遵循2011年奥赛的一门八讲课程；仅有的例外是\autoref{sec:1}与\hyperref[sec:5]{第5章}各用了约一讲半，后续章节在2011年则缩短为半讲. 距离这八“讲”的主题太远的内容，不放入各章主体. 另一方面，每章末尾都有讨论部分，介绍推广、联系、旁支主题及不同观点. 在某些情形，例如拥挤运输与分支运输，这些讨论部分相当于相关运输模型的小型综述. 其风格较为非正式：有时证明陈述，有时只给出草图或提及结果.

为了扩大读者群，让尽可能多的人能够阅读，我决定详细解释一些原本可以假定读者已知的记号——不过，回顾它们总是更好. 各章通过称为“回顾”（Memo）或“重要概念”（Important Notion）的专门知识框回顾部分概念\footnote{两者的区别只是：在后一种情形，我本来更愿意将其视为“回顾”，但学生通常不同意.}. 第三种知识框称为“值得了解”（Good to Know!），提供通常不属于非专门方向数学研究生背景知识的额外概念. 随着全书推进，知识框和说明图的密度都会降低.

为求简便，我决定不反复强调至少一个重要问题：可测性. 你可以相信，书中引入的集合、函数和映射都具有所需的可测性，但我不会逐一明确说明. 实际上，这是我对不够严密所作的唯一让步：全书证明是严格的，至少我希望如此，纯数学家或应用数学家都可以将其作为有关主题的参考. 唯一叙述略显非正式的是关于数值方法的\hyperref[sec:6]{第6章}，这里的意思是我们不提供收敛证明或精确的实现细节.

最后但并非最不重要的是：本书有一章数值方法！尤其介绍与偏微分方程联系最紧密的连续方法，而组合性和离散性最强的方法则在讨论部分简述.

\subsection*{这本书写给谁？}
本书采用应用数学家的视角，应用数学家理应是其自然读者. 不过，我希望与更广泛的读者交流. 纯数学家当然也受到欢迎——纯数学与应用数学之间的区别是否有意义，本来就是个人判断. 他们会看到严格的证明，有时受到不同视角的启发；也可能有兴趣发现最优运输能够用在何处、怎样使用，并作出自己的贡献.

更一般地，可以把区分移到“使用最优运输工作的人”和“研究最优运输本身的人”之间. 前者是本书的自然读者，但后者也可能对内容感兴趣. 反过来，本书是否也面向非数学家，例如物理学家、工程师和理论经济学家？这涉及读者需要多少数学背景. 显然，如果他们有足够的数学背景，而且工作领域与书中应用足够接近，了解这些应用背后的内容会很有意思.

需要多少数学也同样是学生面对的问题. 这是一本数学研究生教材. 即使我试图提供工具，帮助回顾所需背景，要充分利用本书，仍然需要一定先修知识. 主要先修内容是测度论与泛函分析. 我特意没有采纳某位“匿名”审稿人关于加入测度论附录的建议. 我的想法是，希望接近最优运输的数学工作者，应当已经对这些主题有所了解：什么是可测函数、什么是测度、主要收敛定理的条件是什么，以及$L^p$与$W^{1,p}$函数、弱收敛等. 回顾框的作用是帮助读者不至于迷失. 对于阅读本书的非数学家，我希望更具体的处理方式能帮助他们看出哪些性质重要而合理；但这些读者也应掌握一些数学语言，而且肯定需要付出额外努力，才能充分受益.

关于读者，数值部分，即\hyperref[sec:6]{第6章}，还值得多说几句. 详细比较各种方法的优缺点、讨论实现中最巧妙的技巧，以及最新算法，都超出本书范围. 因而，已经研究这一领域的人大概用不上这一章. 另一方面，使用最优运输但不了解数值方法的人，或尚未接触最优运输的数值分析工作者，可能会感兴趣.

我也确信，借助其中许多回顾主要背景概念的知识框和书末习题，本书的一部分材料可以用于有关主题的研究生课程.

\subsection*{这本书包含什么？}
在这篇前言和一篇简短的最优运输导论之后——导论主要介绍问题、历史，以及它与数学其他部分的主要联系——本书包含八章. 最重要的两章是\autoref{sec:1}与\hyperref[sec:5]{第5章}，它们构成最优运输的一般理论. 全书前半部分研究把一个给定测度的质量以最优方式运到另一个测度的问题，先用Monge--Kantorovich框架，再采用最小流观点，其中\autoref{sec:1}最重要. 后半部分转向测度会变化的情形，即\hyperref[sec:5]{第5章}以及之后的\hyperref[sec:7]{第7章}和\hyperref[sec:8]{第8章}. \hyperref[sec:6]{第6章}放在第5章之后，是因为Benamou--Brenier方法与Wasserstein空间中的测地线有关.

\autoref{sec:1}介绍Kantorovich对原始Monge问题所作的松弛，以及相关对偶问题，包括Kantorovich势、$c$-循环单调性等. 这些工具用于给出第一个最优映射存在性定理，即Brenier定理. 讨论部分也主要从Kantorovich解释与对偶性出发.

\hyperref[sec:2]{第2章}集中于一维情形，它更简单，也已有许多推论. 随后介绍Knothe映射：这是由一维构件建立的运输映射，并讨论其退化的最优性. 核心概念是单调运输. 讨论部分将一维和单调映射用于数学，例如等周不等式，以及数学以外的领域，例如图像处理中的直方图均衡化.

\hyperref[sec:3]{第3章}讨论第1章未涉及的极限情形. 第1章的结果告诉我们，怎样处理$p\in(1,+\infty)$时形如$|x-y|^p$的运输成本，却不能处理$p=1$这一Monge原本提出的问题. 为此需要额外想法，尤其是通过第二个变分问题选择某个特定极小元. 另一极限情形$p=\infty$也需要类似技术，本章同样详细讨论. 讨论部分介绍处理一般形如$h(y-x)$的凸成本的主要困难与方法，其中不要求严格凸，并允许取无穷值；这在近年是一个活跃课题. 随后考虑$0<p<1$，即关于距离为凹函数的成本.

\hyperref[sec:4]{第4章}介绍更具欧拉风格的其他表述：怎样通过具有指定散度的向量场$w$描述运输现象，并通过涉及$w$的泛函使总成本最小. 最小化$w$的$L^1$范数时，得到与Monge原始问题等价的表述. 本章主体给出在拉格朗日与欧拉框架之间往返的对应，并研究这一极小化问题及其解. 讨论部分提出两个变体：交通拥挤，对应关于$w$严格凸的成本；分支运输，对应关于$w$为凹的成本.

\hyperref[sec:5]{第5章}介绍最优运输理论的另一个基本工具：运输成本在测度空间上诱导的距离，称为Wasserstein距离. 研究其性质之后，我们研究这些空间中的曲线，尤其是测地线，并强调其与连续性方程的联系. 讨论部分比较Wasserstein距离与概率测度上的其他距离，最后描述测度重心的应用.

\hyperref[sec:6]{第6章}从前一章的想法出发，提出数值方法. 在描述Wasserstein空间时，可以看到求最优运输等价于在连续性方程的解中最小化动能泛函. 这提供了第一个最优运输数值方法，即Benamou--Brenier方法. 本章其余部分描述另两种连续数值方法，讨论部分则介绍离散和半离散方法.

\hyperref[sec:7]{第7章}汇集定义在测度上的有用泛函，并研究其性质，不一定都与最优运输有关，例如凸性、半连续性以及第一变分的计算. 这些泛函经常在建模问题中与运输距离一起出现. 本章还详细介绍位移凸性，即沿Wasserstein空间测地线的凸性. 讨论部分比较多样，既有几何不等式应用，也有城市区域中的均衡问题.

\hyperref[sec:8]{第8章}以独特的方式介绍最优运输最引人注目的应用之一：Wasserstein空间中的梯度流，它使我们能够处理许多演化方程，尤其是抛物型方程. 我们介绍一般框架，并详细研究Fokker--Planck方程. 讨论部分介绍具有这种梯度流结构的其他方程，以及最优运输虽有作用、却并非梯度流的其他演化方程.

在书末详细的参考文献与索引之前，还有来自各章、难度各异的69道习题. 从学生到资深研究者，我们都邀请读者动手尝试这些习题，体会最优运输的滋味.

\begin{flushright}法国奥赛\hfill Filippo Santambrogio\\2015年5月\end{flushright}
